\documentclass[10pt,a4paper]{amsart}
\usepackage[margin=19mm]{geometry}
\usepackage{amsmath,amssymb,mathtools}
\usepackage[scr=rsfs,cal=euler]{mathalfa}
\usepackage{xfrac}
\usepackage[unicode,hidelinks,bookmarks=false]{hyperref}
\newcommand{\ie}{i.e.\ }
\newcommand{\cf}{cf.\ }
\newcommand{\resp}{respectively\ }
\newcommand{\Subsection}{\subsection}
\providecommand{\nobreakdash}{\nobreak\hbox{-}}
\setlength{\emergencystretch}{2em}
\multlinegap0pt
\usepackage{amssymb}
\usepackage{mathtools}
\newtheorem{lemma}{Lemma}
\newcommand{\C}{\mathbb{C}}
\newcommand{\id}{\operatorname{Id}}
\renewcommand{\le}{\leqslant}
\newcommand{\norm}[1]{\left\lVert #1\right\rVert}
\title{There is no universal separable Banach algebra}
\author{Tomasz Kania}
\date{Source: arXiv:2511.04314v3 (CC BY 4.0)}
\begin{document}
\maketitle

\noindent\emph{Source and licence.} There is no universal separable Banach algebra. Version arXiv:2511.04314v3 (CC BY 4.0), distributed by arXiv under the Creative Commons Attribution 4.0 International licence, http://creativecommons.org/licenses/by/4.0/, as recorded in the arXiv licence metadata for this exact version and shown on the abstract page https://arxiv.org/abs/2511.04314v3. Full licence notice: \texttt{licenses/arxiv-2511.04314-CC-BY-4.0.txt}.\par\smallskip

\noindent\textit{Editorial scope.} Counted unit: the article's lemma lem:factor-id (source0.tex lines 177--183). The article's complete original proof of it is delivered with it (source0.tex lines 185--232, one \texttt{proof} environment of 48 lines). No mathematics is written, repaired or re-derived: the statement and the proof are the article's own lines, in the article's own notation, and no retained line is substituted or re-flowed.  Retained uncounted as the source-local context the proof needs: lines 155--158.  Nothing was omitted from any retained span, so the package carries no omission record.  No retained line contains a citation, so the wrapper carries no bibliography and no bibliography entry.  No retained line contains a figure environment, an image-inclusion command or drawing code, so no image is delivered and the package declares no asset dependency.  Editorial formatting only, in addition: an AMS \texttt{amsart} preamble that restates the article's own declarations and macros used by the retained lines, namely the environments \texttt{lemma} on separate counters, together with the standard packages those lines need.  The retained environments print in this document as Lemma 1 in order of appearance, whereas the article prints the same items with the numbering of its own full document.  The complete native source of the article as distributed in the arXiv e-print for this exact version is bundled unchanged as \texttt{sources/188577/source0.tex}.
\begin{equation}\label{eq:product}
(f,x,\alpha)(g,y,\beta) := \bigl(0,0,f(y)+g(x)\bigr)
\qquad (f,g\in X^*,\ x,y\in X,\ \alpha,\beta\in \C).
\end{equation}

\begin{lemma}\label{lem:factor-id}
Let $X$ be a closed subspace of $c_0$, let $B$ be a Banach algebra, and let
\[
\phi:A_X\longrightarrow B
\]
be an injective homomorphism. Then the identity operator on $X^*$ factors through $B$.
\end{lemma}

\begin{proof}
Define bounded linear maps
\[
p:X^*\to B,\qquad p(f):=\phi(f,0,0),
\]
\[
q:X\to B,\qquad q(x):=\phi(0,x,0),
\]
and
\[
j:\C\to B,\qquad j(\alpha):=\phi(0,0,\alpha).
\]
Since $\phi$ is injective, $j$ is injective, and therefore $e:=j(1)\neq 0$.

For $f\in X^*$ and $x\in X$, using \eqref{eq:product} we get
\begin{equation}\label{eq:pxqx}
p(f)q(x)
= \phi(f,0,0)\,\phi(0,x,0)
= \phi\bigl((f,0,0)(0,x,0)\bigr)
= \phi(0,0,f(x))
= f(x)e.
\end{equation}

Choose $\lambda\in B^*$ such that $\lambda(e)=1$ by the Hahn--Banach theorem. Define
\[
S:B\to X^*,
\qquad
(Sb)(x):=\lambda\bigl(b\,q(x)\bigr)
\qquad (b\in B,\ x\in X).
\]
This is a bounded linear operator because
\[
|(Sb)(x)| \le \norm{\lambda}\,\norm{b}\,\norm{q(x)}
\le \norm{\lambda}\,\norm{q}\,\norm{b}\,\norm{x}.
\]
Now \eqref{eq:pxqx} implies that for every $f\in X^*$ and $x\in X$,
\[
\bigl(Sp(f)\bigr)(x)
= \lambda\bigl(p(f)q(x)\bigr)
= \lambda\bigl(f(x)e\bigr)
= f(x).
\]
Hence $Sp(f)=f$ for all $f\in X^*$, that is,
\[
S\circ p = \id_{X^*}.
\]
Thus $\id_{X^*}$ factors through $B$. In particular, $pS$ is a bounded projection from $B$ onto $p(X^*)$, and $p(X^*)$ is a complemented subspace of $B$ isomorphic to $X^*$.
\end{proof}

\end{document}
